% 사물인터넷과 빅데이터 - 1주차 강의개요
% XeLaTeX로 컴파일: xelatex 파일명.tex

\documentclass[11pt, a4paper]{article}

% 한국어 지원 (fontspec만 사용)
\usepackage{fontspec}
\setmainfont{Noto Sans CJK KR}
\setsansfont{Noto Sans CJK KR}

% 페이지 설정
\usepackage[top=2.5cm, bottom=2.5cm, left=2.5cm, right=2.5cm]{geometry}
\usepackage{setspace}
\onehalfspacing

% 색상 및 박스
\usepackage[dvipsnames]{xcolor}
\usepackage{tcolorbox}
\tcbuselibrary{skins, breakable}

% 표 및 그래픽
\usepackage{booktabs}
\usepackage{tabularx}
\usepackage{graphicx}
\usepackage{float}
\usepackage{amssymb}

% 목록 및 링크
\usepackage{enumitem}
\usepackage{hyperref}
\hypersetup{
    colorlinks=true,
    linkcolor=NavyBlue,
    urlcolor=RoyalBlue
}

% 헤더/푸터
\usepackage{fancyhdr}
\pagestyle{fancy}
\fancyhf{}
\fancyhead[L]{\small\textcolor{gray}{사물인터넷과 빅데이터}}
\fancyhead[R]{\small\textcolor{gray}{1주차: 강의개요}}
\fancyfoot[C]{\thepage}
\renewcommand{\headrulewidth}{0.4pt}

% 섹션 스타일
\usepackage{titlesec}
\titleformat{\section}
    {\Large\bfseries\color{NavyBlue}}
    {\thesection}{1em}{}[\titlerule]
\titleformat{\subsection}
    {\large\bfseries\color{RoyalBlue}}
    {\thesubsection}{1em}{}

% 커스텀 박스 정의
\newtcolorbox{keypoint}{
    colback=blue!5,
    colframe=NavyBlue,
    fonttitle=\bfseries,
    title=핵심 메시지,
    breakable
}

\newtcolorbox{activity}{
    colback=green!5,
    colframe=ForestGreen,
    fonttitle=\bfseries,
    title=활동,
    breakable
}

\newtcolorbox{info}{
    colback=gray!5,
    colframe=gray,
    breakable
}

% 목차/표/그림 한글화
\renewcommand{\contentsname}{목차}
\renewcommand{\tablename}{표}
\renewcommand{\figurename}{그림}

\begin{document}

%========================================
% 표지
%========================================
\begin{titlepage}
\centering
\vspace*{3cm}

{\Huge\bfseries\color{NavyBlue} 사물인터넷과 빅데이터}\\[1cm]
{\LARGE 1주차 강의개요}\\[0.5cm]
{\Large 오리엔테이션 및 강좌 소개}\\[2cm]

\begin{info}
\centering
\begin{tabular}{l l}
\textbf{문서 버전} & v1.0.0.0 \\
\textbf{작성일} & 2026-01-28 \\
\textbf{작성자} & Claude \\
\end{tabular}
\end{info}

\vfill

{\large 청주교육대학교 컴퓨터교육과}\\[0.5cm]
{\large 양창모}

\end{titlepage}

%========================================
\section{차시 정보}
%========================================

\begin{table}[H]
\centering
\begin{tabularx}{\textwidth}{lX}
\toprule
\textbf{항목} & \textbf{내용} \\
\midrule
주차/차시 & 1주차 (단일 차시) \\
주제 & 오리엔테이션 및 강좌 소개 \\
수업 시간 & 100분 (50분 × 2) \\
운영 방식 & 대면 수업 \\
바이브코딩 & $\checkmark$ 활용 \\
\bottomrule
\end{tabularx}
\caption{1주차 차시 정보}
\end{table}

%========================================
\section{학습 목표}
%========================================

본 차시를 마치면 학생은 다음을 할 수 있다.

\begin{enumerate}
    \item 강좌의 목표와 진행 방식을 설명할 수 있다.
    \item 바이브코딩의 개념과 가능성을 체험을 통해 이해할 수 있다.
    \item 마이크로비트 LED에 자신의 이름을 표시하는 프로그램을 바이브코딩으로 만들 수 있다.
\end{enumerate}

%========================================
\section{핵심 메시지}
%========================================

\begin{keypoint}
\begin{enumerate}
    \item ``데이터는 원인을 말하지 않는다. 우리가 원인을 '추정'할 뿐이다.''
    \item ``AI는 집중도를 판단하지 않는다. 우리가 만든 기준을 계산할 뿐이다.''
    \item ``이 수업은 코딩을 가르치지 않는다. 코딩을 의심하는 방법을 가르친다.''
\end{enumerate}
\end{keypoint}

%========================================
\section{수업 흐름}
%========================================

\begin{table}[H]
\centering
\begin{tabularx}{\textwidth}{c l X l}
\toprule
\textbf{시간} & \textbf{활동} & \textbf{내용} & \textbf{자료/도구} \\
\midrule
20분 & 강좌 소개 & 목표, 15주 흐름, 평가 방식 안내 & 강의슬라이드 \\
30분 & 바이브코딩 시연 & 수업집중도분석프로그램 5개 탭 시연 & 시뮬레이터 \\
10분 & 예시 프로젝트 & ``수업 집중도'' 프로젝트 전체 흐름 & 강의슬라이드 \\
30분 & LED 실습 & 마이크로비트 LED에 이름 표시 & MakeCode, Claude \\
10분 & 마무리 & 2주차 예고, 과제 안내 & 강의슬라이드 \\
\bottomrule
\end{tabularx}
\caption{100분 수업 흐름}
\end{table}

%========================================
\section{주요 내용}
%========================================

\subsection{강좌 소개 (20분)}

\textbf{목표}: 강좌 전체 흐름과 평가 방식 이해

\begin{itemize}
    \item 강좌명: 사물인터넷과 빅데이터
    \item 핵심 흐름: 센서 → 데이터 수집 → 시각화 → 분석 → AI 모델
    \item 평가: 주간과제 80\% + 최종프로젝트 15\% + 성찰보고서 5\%
    \item 2022 개정 교육과정 연계 (피지컬 컴퓨팅, 데이터 수집·분석)
\end{itemize}

\subsection{바이브코딩 프로그램 시연 (30분)}

\textbf{목표}: 바이브코딩의 가능성을 직접 확인

\textbf{수업집중도분석프로그램} 5개 탭 시연:

\begin{enumerate}
    \item \textbf{탭 1 - 데이터 입력}: 온도, 소음, 집중도 데이터 입력
    \item \textbf{탭 2 - 데이터 분석}: 산점도, 상관계수, 패턴 분석
    \item \textbf{탭 3 - 학생 모델}: 규칙 기반 AI (IF-THEN)
    \item \textbf{탭 4 - 교수 모델}: 머신러닝 AI
    \item \textbf{탭 5 - 모델 비교}: 정확도 비교, 해석 가능성 논의
\end{enumerate}

\begin{info}
\textbf{핵심 질문}: ``이 프로그램, 코드를 한 줄도 직접 작성하지 않고 만들었어요!''
\end{info}

\subsection{예시 프로젝트 소개 (10분)}

\textbf{목표}: 15주 프로젝트 흐름 이해

``수업 집중도'' 프로젝트 단계: 기획 → 제작 → 수집 → 분석 → 발표

\begin{info}
\textbf{핵심 질문}: ``온도가 높으면 집중도가 낮아질까?''\\
→ 상관관계 vs 인과관계 구분 필요성 강조
\end{info}

\subsection{마이크로비트 LED 실습 (30분)}

\textbf{목표}: 바이브코딩 첫 체험

\begin{activity}
\textbf{실습 단계}:
\begin{enumerate}
    \item Claude에게 MakeCode 블록 코드 요청
    \item MakeCode에서 프로그램 작성
    \item 시뮬레이터에서 결과 확인
    \item 실제 마이크로비트에서 실행
    \item 실행 모습 사진 촬영
\end{enumerate}
\end{activity}

\textbf{프롬프트 예시}:
\begin{quote}
``마이크로비트 LED에 '홍길동'을 한 글자씩 순서대로 표시하는 MakeCode 블록 코드를 알려줘. 각 글자는 1초간 표시되고 다음 글자로 넘어가도록 해줘.''
\end{quote}

\subsection{마무리 및 과제 안내 (10분)}

\textbf{2주차 예고}:
\begin{itemize}
    \item 바이브코딩 개념 심화
    \item Claude 사용법
    \item 센서 탐색 (내장 + Grove)
\end{itemize}

\textbf{1주차 과제}:
\begin{itemize}
    \item 제출물: 실행 사진 + 대화목록 전문(PDF) + 대화목록 요약본(PDF)
    \item 기한: 일요일 자정 (LMS)
    \item 특이사항: 평가 점검용 (모든 학생 10\% 만점 부여)
\end{itemize}

%========================================
\section{교수-학생 활동}
%========================================

\begin{table}[H]
\centering
\begin{tabularx}{\textwidth}{l X}
\toprule
\textbf{교수 활동} & \textbf{내용} \\
\midrule
강좌 소개 & 목표, 진행 방식, 평가 안내 \\
시연 & 바이브코딩 결과물(수업집중도분석프로그램) 시연 \\
안내 & 예시 프로젝트 흐름, LED 실습 가이드 \\
지원 & 개별 실습 피드백 \\
\bottomrule
\end{tabularx}
\caption{교수 활동}
\end{table}

\begin{table}[H]
\centering
\begin{tabularx}{\textwidth}{l X}
\toprule
\textbf{학생 활동} & \textbf{내용} \\
\midrule
체험 & 바이브코딩 결과물 관찰 \\
실습 & LED 이름 표시 프로그램 제작 (바이브코딩) \\
제출 & 사진 + 대화목록 PDF \\
\bottomrule
\end{tabularx}
\caption{학생 활동}
\end{table}

%========================================
\section{평가 연계}
%========================================

\begin{table}[H]
\centering
\begin{tabularx}{\textwidth}{l X}
\toprule
\textbf{항목} & \textbf{내용} \\
\midrule
과제 & LED 이름 표시 \\
배점 & 10\% (점검용 만점 처리) \\
제출물 & 사진 + 대화목록 전문(PDF) + 대화목록 요약본(PDF) \\
기한 & 일요일 자정 \\
\bottomrule
\end{tabularx}
\caption{1주차 평가 연계}
\end{table}

%========================================
\section{필요 자료 및 도구}
%========================================

\subsection{시뮬레이터/프로그램}

\begin{table}[H]
\centering
\begin{tabularx}{\textwidth}{c l X}
\toprule
\textbf{\#} & \textbf{자료명} & \textbf{용도} \\
\midrule
1 & 수업집중도분석프로그램\_v2.1.0.0 & 바이브코딩 시연 (30분) \\
2 & MakeCode시뮬레이터-LED이름표시 & LED 실습 참고 \\
\bottomrule
\end{tabularx}
\caption{시뮬레이터/프로그램}
\end{table}

\subsection{하드웨어}

\begin{table}[H]
\centering
\begin{tabularx}{\textwidth}{c l c X}
\toprule
\textbf{\#} & \textbf{항목} & \textbf{수량} & \textbf{비고} \\
\midrule
1 & 마이크로비트 V2 & 21개 & 학생 1인 1개 \\
2 & USB 케이블 & 21개 & 데이터 전송용 \\
\bottomrule
\end{tabularx}
\caption{하드웨어}
\end{table}

\subsection{소프트웨어}

\begin{table}[H]
\centering
\begin{tabularx}{\textwidth}{c l l X}
\toprule
\textbf{\#} & \textbf{도구} & \textbf{URL} & \textbf{용도} \\
\midrule
1 & MakeCode & makecode.microbit.org & 블록 코딩 \\
2 & Claude & claude.ai & 바이브코딩 \\
\bottomrule
\end{tabularx}
\caption{소프트웨어}
\end{table}

%========================================
\section{그림 목록}
%========================================

\begin{table}[H]
\centering
\begin{tabularx}{\textwidth}{c X l l}
\toprule
\textbf{\#} & \textbf{그림 제목} & \textbf{용도} & \textbf{형식} \\
\midrule
1 & 15주 전체 흐름도 & 강좌 소개 & TikZ \\
2 & 시뮬레이터 탭 1-5 화면 & 바이브코딩 시연 & TikZ (5개) \\
3 & 바이브코딩 vs 전통 프로그래밍 & 개념 설명 & TikZ \\
4 & 마이크로비트 V2 구성도 & 하드웨어 소개 & 이미지 \\
5 & MakeCode 화면 구성 & 실습 안내 & TikZ \\
6 & LED 이름 표시 실행 화면 & 실습 결과 & TikZ \\
7 & Grove 센서 이미지 갤러리 & 센서 소개 & 이미지 (10종) \\
\bottomrule
\end{tabularx}
\caption{그림 목록}
\end{table}

%========================================
\section{임용 연결}
%========================================

\begin{table}[H]
\centering
\begin{tabularx}{\textwidth}{l X}
\toprule
\textbf{연결 포인트} & \textbf{내용} \\
\midrule
2022 개정 교육과정 & ``피지컬 컴퓨팅이 초등 실과에 들어갔어요'' \\
데이터 리터러시 & ``데이터 수집·분석이 교육과정에 포함되었어요'' \\
\bottomrule
\end{tabularx}
\caption{임용 연결 포인트}
\end{table}

%========================================
\section*{참고자료}
%========================================

\begin{table}[H]
\centering
\begin{tabularx}{\textwidth}{c l l X}
\toprule
\textbf{\#} & \textbf{자료명} & \textbf{유형} & \textbf{비고} \\
\midrule
1 & 강의계획서\_사물인터넷과빅데이터\_v1.6.0.0 & md & 1주차 내용 참조 \\
2 & 수업집중도분석프로그램\_v2.1.0.0 & HTML & 시연용 프로그램 \\
3 & MakeCode시뮬레이터-LED이름표시 & HTML & 실습용 시뮬레이터 \\
\bottomrule
\end{tabularx}
\end{table}

%========================================
\section*{생성/수정 이력}
%========================================

\begin{table}[H]
\centering
\begin{tabularx}{\textwidth}{l l l l X l}
\toprule
\textbf{버전} & \textbf{날짜} & \textbf{시간} & \textbf{변경 수준} & \textbf{변경 내용} & \textbf{작성자} \\
\midrule
v1.0.0.0 & 2026-01-28 & 05:20 & 최초 & 초안 작성 & Claude \\
\bottomrule
\end{tabularx}
\end{table}

\end{document}
